\chapter{Database}

Il database è il cuore della persistenza dati di Ruggine. È progettato per garantire integrità, coerenza e performance, supportando autenticazione, chat, gruppi e inviti. La progettazione segue principi di normalizzazione e sicurezza.

\section{Struttura e tabelle principali}
Le principali tabelle del database sono: \texttt{users}, che contiene gli utenti registrati (id, username, password hash); \texttt{groups}, che rappresenta gruppi e chat private (id, name, is\_dm); \texttt{group\_members}, tabella ponte N-a-N tra utenti e gruppi; \texttt{messages}, che memorizza i messaggi (id, group\_id, sender\_user\_id, content, created\_at); \texttt{invitations}, che gestisce inviti a gruppi (id, from\_user\_id, to\_user\_id, group\_id, status, timestamps).

\subsection{Esempio: creazione tabella messaggi}
\begin{lstlisting}[language=SQL, caption={Estratto da migration SQL}]
CREATE TABLE IF NOT EXISTS messages (
	id INTEGER PRIMARY KEY AUTOINCREMENT,
	group_id INTEGER NOT NULL,
	sender_user_id INTEGER NOT NULL,
	content TEXT NOT NULL,
	created_at DATETIME DEFAULT (datetime('now')),
	FOREIGN KEY (group_id) REFERENCES groups(id) ON DELETE CASCADE,
	FOREIGN KEY (sender_user_id) REFERENCES users(id) ON DELETE CASCADE
);
\end{lstlisting}


\section{Flussi tipici ed esempi SQL}
\subsection{Registrazione utente}
Quando un nuovo utente si registra, viene inserito nella tabella \texttt{users} e aggiunto al gruppo globale tramite la tabella ponte \texttt{group\_members}.
\begin{lstlisting}[language=SQL, caption={Inserimento nuovo utente}]
INSERT INTO users (username, password_hash) VALUES ('alice', 'hashpassword');
INSERT INTO group_members (group_id, user_id) VALUES (1, last_insert_rowid());
\end{lstlisting}

\subsection{Invio messaggio}
Quando un utente invia un messaggio in un gruppo, viene inserita una riga nella tabella \texttt{messages} con riferimenti a gruppo e utente.
\begin{lstlisting}[language=SQL, caption={Inserimento messaggio in chat}]
INSERT INTO messages (group_id, sender_user_id, content) VALUES (2, 1, 'Ciao a tutti!');
\end{lstlisting}

\subsection{Gestione inviti}
La creazione di un invito comporta l'inserimento nella tabella \texttt{invitations}. L'accettazione o il rifiuto aggiorna lo stato e il trigger aggiorna il timestamp.
\begin{lstlisting}[language=SQL, caption={Gestione inviti}]
INSERT INTO invitations (from_user_id, to_user_id, group_id) VALUES (1, 2, 3);
UPDATE invitations SET status = 'accepted' WHERE id = 1;
-- Il trigger aggiorna updated_at automaticamente
\end{lstlisting}

\subsection{Esempio di query: recupero messaggi di un gruppo}
\begin{lstlisting}[language=SQL, caption={Recupero messaggi di un gruppo}]
SELECT m.content, u.username, m.created_at
FROM messages m
JOIN users u ON m.sender_user_id = u.id
WHERE m.group_id = 2
ORDER BY m.created_at ASC;
\end{lstlisting}

\section{Tecnologie e motivazioni}
Il database scelto è SQLite per semplicità e portabilità in fase di sviluppo, ma la struttura è compatibile con PostgreSQL per ambienti di produzione. L'accesso avviene tramite SQLx, che garantisce query asincrone e sicure.

\section{Sicurezza e integrità}
\begin{itemize}
	\item Vincoli di chiave esterna e ON DELETE CASCADE per coerenza referenziale.
	\item Query parametrizzate per prevenire SQL injection.
	\item Migrazioni versionate per tracciare ogni modifica allo schema.
	\item Trigger per aggiornare automaticamente i timestamp.
\end{itemize}

\section{Considerazioni progettuali}
La normalizzazione facilita l'estensione futura (es. ruoli, permessi). La scelta di SQLx e delle migrazioni versionate garantisce riproducibilità e sicurezza. L'architettura a tabelle separate isola le responsabilità e semplifica la manutenzione.